\section{Introduction}
\label{sec:intro}

A launchpad is a platform that lets users create their own cryptocurrencies. On pump.fun, the largest launchpad on Solana, any user can create a token by choosing a name, a ticker and an image, with no need to write or deploy a smart contract. Tokens created on pump.fun trade on a bonding curve, which sets the price from the first trade onward. As a result, these launchpads produce an enormous number of small, new cryptocurrencies. One study of pump.fun counts about 15 million coins launched over two years~\citep{szwajcok2026factories}, while another finds that fewer than 2\% of the platform's tokens reach a major decentralized exchange~\citep{mancino2025memecoin}. In the public corpus we use, about 0.7\% of launched tokens graduated from the bonding curve during our 39-day window~\citep{slinkdev2026corpus}.

For most of these tokens, trading dries up after a handful of transactions and they are forgotten. For a small subset, however, the price surges: within a few seconds, strong buying pressure drives a steep, abrupt rise, which typically fades within minutes. In a stratified sample of tokens from our data, the median surge peaked 73\% above its starting price about 32 seconds after it began. Such surges attract retail traders and automated trading bots, and some are the product of coordinated manipulation~\citep{szwajcok2026factories,mongardini2025midsummer}. Understanding and predicting them therefore matters to traders, to platforms and wallet applications deciding whether to warn users, and to researchers and regulators.

Prior work has studied this market from several angles. Measurement studies describe the platform, its graduation dynamics and the manipulation strategies used on it~\citep{mancino2025memecoin,kamat2026auditing,szwajcok2026factories}. Predictive studies estimate whether a token will graduate from its bonding curve~\citep{marino2026predicting}, flag likely rug pulls from the first minutes of trading~\citep{li2026catching}, or identify high-risk launches before they migrate to an exchange~\citep{hu2026melt}. A separate line of work studies pump-and-dump schemes on centralized exchanges, which are typically organized in advance on platforms such as Telegram~\citep{kamps2018moon,xu2019anatomy,hu2023sequence}. Detection systems spot such pumps seconds after they start~\citep{lamorgia2020pump,lamorgia2023doge}, and recent work ranks likely target coins shortly before an announced pump~\citep{bolz2024realtime}. Our work differs from all of these in one important way: we ask whether a surge is about to begin, within seconds, at any arbitrary point in a token's life, using only on-chain data. To our knowledge, this prediction task has not been studied before. \citet{marino2026predicting}, for instance, predict graduation and explicitly leave prediction of the pump itself to future work.

This question is easy to get wrong. Labels depend on the future by construction, features computed from on-chain data can quietly include future information, and the market changes from week to week. We therefore put as much effort into the evaluation protocol as into the models. All model and hyperparameter choices were made on walk-forward folds inside the training period. A separate validation week was used for development reporting. The final week was locked until we had written down, and committed to version control, exactly which models and metrics we would report. We then scored it once.

Our contributions are as follows.
\begin{enumerate}
  \item \textbf{A task and benchmark.} We formulate seconds-ahead surge-onset prediction on a bonding-curve launchpad, and we build a point-in-time dataset of 8.4 million observations of 622,303 pump.fun tokens with a prespecified surge detector, causal observation sampling and chronological, token-disjoint splits.
  \item \textbf{A preregistered evaluation.} We evaluate five model families on a locked final week, scored once under a protocol committed before the test data were read. We report token-clustered confidence intervals, lead times, event-level recall and results within strata of token age and activity.
  \item \textbf{A predictability result.} Surge onset is predictable from causal on-chain state. On the held-out week, the best model reaches an ROC-AUC of 0.866 at a 60-second horizon and 0.768 among actively traded tokens, and its top-ranked observations contain surge onsets at about six times the base rate.
  \item \textbf{A model-family comparison.} Gradient-boosted trees and L2-regularized logistic regression perform almost identically on the same engineered features, and neither wins in every subgroup. A multimodal Transformer trained on raw event sequences, with self-supervised pretraining, performs worse than both overall.
  \item \textbf{A leakage case study.} We show how a reasonable-looking, cost-driven choice of which transactions to fetch from an RPC node leaked future information and inflated a model's ROC-AUC from 0.835 to 0.917, even though every individual record passed timestamp checks. We also describe a separate evaluation bug and the tests that caught both problems.
\end{enumerate}

We also find that predictability and profitability are different things. The observations our model ranks highest are, on average, followed by steep losses, which is consistent with a market where any easily detected edge is short-lived. We do not offer a trading strategy, and we discuss the dual-use risks of this kind of model in Section~\ref{sec:discussion}.

Section~\ref{sec:background} explains how pump.fun works and Section~\ref{sec:related} reviews related work. Sections~\ref{sec:problem}--\ref{sec:protocol} define the task, the data, the models and the evaluation protocol. Section~\ref{sec:results} reports results, Section~\ref{sec:leakage} presents the leakage case study, Section~\ref{sec:discussion} discusses implications, limitations and ethics, and Section~\ref{sec:conclusion} concludes.
